\zhbox{54}{64}{235}{655}{10}{\textbf{电源}\par
MP2984 的内部电路由 AVDD（5.2 V）供电，栅极驱动器由 VCC（7.2 V）供电。VIN 和 VOUT 调节 VCC 电压（$V_{CC}$），VCC 为 AVDD 电压（$V_{AVDD}$）供电。\par
若 VIN 已供电且 $V_{EN}$ 为高，MP2984 将 $V_{CC}$ 调节为7.2 V、$V_{AVDD}$ 调节为5.2 V。若 $V_{AVDD}$ 超过 UVLO 阈值且 ENPWR 为高，器件开始开关工作，并通过软启动控制调节 $V_{OUT}$。若 $V_{IN}$ 和 $V_{OUT}$ 均超过8.8 V，VCC 由两者中电压较低的电源供电以降低损耗；否则，由电压较高的电源供电，以确保 VCC 电压充足。VCC 与 BST 引脚各有独立 UVLO 保护，可使栅极信号保持关闭。UVLO 触发后，应确保 VCC 与 BST 电压足以使器件重新启动；这不适用于 AVDD UVLO。\par
MP2984 的输入工作范围为6--36 V。若启动后 VCC 由 VOUT 供电，器件可持续工作直到 $V_{IN}$ 低于5 V。\par
若器件因 AVDD\_UVLO 或 EN 信号而关闭，I\textsuperscript{2}C 接口无法响应主机，COMP 被拉低。$V_{CC}$、$V_{AVDD}$ 与自举电压（$V_{BST}$）随漏电缓慢下降，所有逻辑关闭。\par
\textbf{启动}\par
MP2984 通过软启动（SS）启动 IC。软启动电路向 SS 引脚充电，使软启动电压（$V_{SS}$）从0 V缓慢上升。软启动电流（$I_{SS}$）随后送入误差放大器（EA）以控制输出。$V_{SS}$ 超过由 VREF 位设置的 $V_{REF}$ 后，软启动完成，器件进入闭环调节。稳态时，除非触发保护，$V_{SS}$ 应上升并钳位在0.6 V（高于 $V_{REF}$）。\par
启动后，由于 $V_{OUT}<V_{IN}$，MP2984 通常先进行降压开关工作。若输出存在预偏置，器件应在 $V_{SS}$ 超过 $V_{FB}$ 后才开关工作；$V_{FB}$ 与输出预偏置电压成正比。软启动期间，IC 工作于自动 PFM 模式，过压保护（OVP）和打嗝模式过流保护（OCP）不起作用。}
\zhbox{324}{64}{235}{346}{10}{\textbf{使能（EN）及可配置欠压锁定（UVLO）保护}\par
EN 引脚使能或禁用 MP2984。若 EN 电压（$V_{EN}$）超过高阈值（$>1.1$ V），部分内部电路启动，称为微功耗模式。若 $V_{EN}$ 超过启动阈值（1.35 V），IC 启动并使能所有功能。若 $V_{EN}$ 降至低阈值以下（$<1.28$ V），则禁用开关工作。\par
若 $V_{EN}<0.4$ V，MP2984 关闭。关闭后，器件从输入电源吸取很小的电流（典型值小于1 $\mu$A）。EN 可承受最高40 V电压。将 EN 与 VIN 相连可自动启动。通过 EN 关断时，经过200 ms放电时间（$t_{DISCHARGE}$）后，I\textsuperscript{2}C 接口复位至默认值。\par
MP2984 具有可配置 UVLO 迟滞。启动期间，一旦 $V_{EN}>1.35$ V，EN 引脚就输出4.7 $\mu$A电流（见图6）。IC 启动后，$V_{IN}$ 必须下降以克服该电流源，才能停止开关工作。输入开关开启阈值（见图6）可由式(1)计算：}
\zhbox{324}{452}{234}{31}{10}{输入开关关断阈值可由式(2)计算：}
\zhbox{324}{541}{234}{26}{10}{其中 $V_{EN\_ON}$ 典型值约1.35 V，$V_{EN\_OFF}$ 约1.28 V。}
\zhbox{324}{698}{234}{18}{10}{\centering\textbf{图6：可配置输入 UVLO}}
